% !TEX root = ../main.tex
\lecture{10}{2026-11-16}{Geldflussrechnung}

\section{Geldflussrechnung}

\subsection{Wesen und Bedeutung}
Ist eine zusätzliche Betrachtung der Geldströme/Ein oder
Auszahlungen von Kunden.
\begin{itemize}
    \item Geldflussrechnung GFR oder cash flow statement (CFS)
    \item Fokus auf effektiven geldströme (Liquidität)
    \begin{itemize}
        \item Wenn nicht genug Liquidität vorhanden um
        Verpflichtungen zu erfüllen (Löhne, Schulden, usw.) -> Konkurs
    \end{itemize}
    \item Wenig spielraum (gegensatz Bewertung worin
    Abschreibungszeitraum gewählt werden kann) entweder fliesst
    Geld rein oder raus.
    \item Engerer Zusammenhang zwischen Bilanz und ER (Zeigt
    innerhalb der Periode)
    \item 
\end{itemize}

\subsubsection{Geld}
\begin{definition}
    [Geld/cash] ist ein Fonds (Mittelgesamtheit nicht aus einzelner
    Positionen): Flüssige Mittel (meist Umlaufvermögen oder < 3
    Monate Frist)
    \begin{itemize}
        \item Kasse
        \item Bankkonto
        \item Postkonto
        \item kfr. Festgeldanlagen
    \end{itemize}
    Nettoflüssigen Mittel:
    \begin{itemize}
        \item Ist gleich Flüssige Mittel - kfr.
        Bankverbindlichkeiten (z.B Kontokorrent auf kfr. FK)
    \end{itemize}
\end{definition}


\subsubsection{Beispiel}

\begin{itemize}
    \item Debitoren AB: 1000
    \item Debitoren SB: 600
    \item Warenertrag: 6000
    \item GFR: AB + Zwischenstandveränderung + Veränderung
    AB/SB = SB 
    \begin{itemize}
        \item 1000 + 6000 + Geldfluss = 600
        \item Geldfluss = 6400
    \end{itemize}
\end{itemize}

\subsubsection{Schema}
Wie viel der Veränderung des Debitoren Kontos bestand aus Kundenzahlungen und welche aus Verluste

\textbf{Einzahlungen von Kunden (Geldzufluss)}
\[
GF = Warenertrag +/-(Abnahme/Zunahme) Brutto-Debitoren - Debitorenverluste +/- Zunahme/Abnahme Delkredere 
\]


\textbf{Auszahlungen an Lieferanten (Geldabfluss)}
\[
GF = - Warenaufwand +/- Abnahme/Zunahme Warenlager  +/- Zunahme/Abnahme Kreditoren 
\]


\subsection{Erstellung}

GFR besteht aus 2 Teilen:

\begin{itemize}
    \item Liquiditätsnachweis
    \begin{itemize}
        \item Fondskonten + Veränderung während der Periode
    \end{itemize}
    \item GFR an sich 
    \begin{itemize}
        \item Ursachen der Veränderung des Fonds 
        \item 3 Teile
        \begin{itemize}
            \item CF aus Betriebstätigkeit
            \begin{itemize}
                \item 2 Darstellungen: Direkte/Indirekte
                \item Indirekt beginnt mit Gewinn/Verlust +/-
                Anpassungen um auf GFR zu kommen
            \end{itemize}
            \item CF aus investitionstätigkeit
            \item CF aus Finanzierungstätigkeit
        \end{itemize}
    \end{itemize}
\end{itemize}